\documentclass[11pt]{article}

% ---------------------------------------------------------------------------
% Does Optimizing the Loop Beat Optimizing the Prompt?
% Pre-print draft. Build: pdflatex main && bibtex main && pdflatex main x2
% Author list is a placeholder (TODO) — fill before submission.
% ---------------------------------------------------------------------------

\usepackage[utf8]{inputenc}
\usepackage[T1]{fontenc}
\usepackage{amsmath,amssymb}
\usepackage{booktabs}
\usepackage{graphicx}
\usepackage{hyperref}
\usepackage{natbib}
\usepackage[margin=1in]{geometry}
\usepackage{xcolor}
\newcommand{\dd}{\ensuremath{\Delta}}
\newcommand{\todo}[1]{\textcolor{red}{[TODO: #1]}}

\title{Your Loop Optimizer Is Not Optimizing Your Loop\\
\large A silent failure mode in componentized prompt optimization,\\
and the one integer that catches it}

\author{Author list TBD \\ \texttt{github.com/loopgepa/whystuck}}
\date{Draft: July 2026}

\begin{document}
\maketitle

\begin{abstract}
Reflective prompt optimizers such as GEPA \citep{agrawal2025gepa} can evolve
\emph{multiple} text components of a compound LLM system, not just a single prompt, and a
natural hypothesis is that optimizing more of a generate--critique--refine loop should beat
optimizing its generation prompt alone. We pre-registered and ran eleven four-arm
experiments testing this on verifiable instruction following \citep{zhou2023ifeval} and on
frontier-hard code \citep{zhuo2024bigcodebench}, and measured $\dd(D-B)\approx 0$
throughout. \textbf{An audit of our own harness then established that the loop-level arm's
treatment had never been applied}, and this paper reports that finding rather than the null
it invalidates. Two defects combined: our reflective dataset was \emph{component-agnostic},
handing the critique gene the \emph{generator's} inputs and outputs, so the optimizer
rewrote a detector while blind to that detector's behaviour --- it duly proposed generation
prose into a critique slot; and GEPA's default selector rewrites one component per
proposal, so the two-gene arm searched the shared productive gene roughly half as hard as
the one-gene arm under a matched cap. Replaying our flagship run, the critique gene is
byte-identical to the seed in $15/15$ winning genomes. We also correct a percentile-index
error that made our one interval-excluding-zero result appear significant: the actuator
run's $\dd(D-B)=-0.10$ has an exact-resample interval of $[-0.150,+0.000]$, which
\emph{includes} zero (sign test $p=0.50$). Both defects are invisible in held-out scores,
which is what makes them worth reporting: any componentized-optimizer study can measure a
confident zero from an experiment that never ran. Findings that survive the audit are the
mechanistic ones: the intrinsic critique is a near-random failure detector
($\mathrm{F}_1\!\approx\!0.11$), consistent with the documented unreliability of LLM
self-verification \citep{huang2023cannot,kamoi2024whencan} and with the independent
observation that detection ability does not predict correction success
\citep{decomposing2026selfcorrection}; \emph{given} a perfect checker at inference, the
repair signal's shape matters ($+0.063$ text, $+0.20$ code, both intervals excluding zero)
while naming the failed constraint alone adds nothing; and the residual localizes to one
mechanical skill (\texttt{exact\_word\_count}, $0.16\!\to\!0.33\!\to\!0.50$). We release
the corrected harness, which reports per-gene accepted-change counts and train aggregates
so that ``the treatment was applied'' is checkable by a reader, together with every
cassette and pre-registration --- including the superseded ones, labelled as such.
\end{abstract}

\section{Introduction}
\label{sec:intro}

Automatic prompt optimization (APO) has matured from single-instruction search
\citep{madaan2023selfrefine} to optimizers that treat an entire LLM program as
the object of optimization. GEPA \citep{agrawal2025gepa}, the current
state-of-the-art reflective optimizer and an ICLR-2026 oral, evolves any set of
text components of a compound system --- prompts, rubrics, tool descriptions ---
by reflecting on execution traces in natural language and maintaining a Pareto
frontier of candidates. DSPy \citep{khattab2023dspy} similarly compiles and
optimizes multi-module pipelines end-to-end.

This capability raises a concrete engineering question that, to our knowledge,
has not been isolated experimentally for a self-refine loop's own verifier. The
adjacent question for multi-\emph{agent} systems has since been benchmarked:
\citet{bai2026masprompt} find that GEPA improves single-agent performance on all five of
their tasks ($+2$ to $+14$ points) while its multi-agent extension ranges from $-16$ to
$+14$ depending on task and workflow topology --- more optimizable surface is not reliably
better. Our question is narrower and about a different object: \emph{given a self-refine loop, does
optimizing the loop's verification/critique component (and beyond) yield better
held-out performance than optimizing only its generation prompt, under a matched
budget?} We call the former \textbf{loop-level} and the latter
\textbf{prompt-only} optimization. The hypothesis is intuitive --- more
optimizable surface, more headroom --- and if true it would justify the extra
cost of evolving critique rubrics and plan modules.

We set out to answer this question with a pre-registered, control-heavy study, measured
$\dd\approx 0$ eleven times, and then found that our own harness had prevented the
loop-level arm from receiving its intervention. What we can contribute is therefore
different from what we planned, and in one respect more useful:
\begin{enumerate}
  \item \textbf{A silent failure mode for componentized optimization}
    (Section~\ref{sec:defect}), established on our own artifact: a
    \emph{component-agnostic reflective dataset} shows every component the generator's
    inputs and outputs, so a critique gene is rewritten by a reflector that cannot see the
    critique; and a one-component-per-proposal selector silently halves a two-gene arm's
    search on the gene that actually moves the metric. Neither is visible in held-out
    scores. Replaying our flagship run, the critique gene is byte-identical to the seed in
    $15/15$ winning genomes and no critique proposal was ever accepted.
  \item \textbf{The diagnostic that catches it}, now part of the released harness: every
    arm reports per-gene accepted-change counts, its train aggregate, and its final gene
    text, so ``the treatment was applied'' is checkable rather than assumed. We recommend
    this as a reporting minimum for any multi-component optimization result.
  \item \textbf{A statistical correction}: our exact-resample interval dropped the largest
    resample atom, which at $n{=}3$ carries $1/27 = 3.7\% > \alpha/2$ and cannot be
    excluded. The one result we had reported as excluding zero does not
    (Section~\ref{sec:defect}).
  \item A taxonomy separating a legitimate non-oracle \emph{measurement tool}
    from oracle contamination (Section~\ref{sec:method}), a distinction a
    pre-spend design review found essential: two seemingly-innocuous statistics
    were the scoring oracle in disguise. This part of the design held up.
  \item A mechanistic account of the loop that survives the audit
    (Section~\ref{sec:analysis}): the intrinsic critique is a near-random detector
    ($\mathrm{F}_1\!\approx\!0.11$); \emph{given} a perfect checker at inference, repair
    feedback shape is worth $+0.063$ (text) and $+0.20$ (code) while naming the failed
    constraint adds nothing; and the residual localizes to a single mechanical skill.
\end{enumerate}

We report every result as-is, including a run in which optimization
\emph{degraded} held-out performance, and including the audit that invalidated our own
headline. The point is not that GEPA is weak --- it is a strong optimizer --- but that
\emph{pointing an optimizer at a component is not the same as optimizing it}, and the
difference does not show up in the scores.

\section{Background and Related Work}
\label{sec:related}

\paragraph{Loop engineering, and which loop this paper is about.}
During June 2026 industry practice converged on a name for designing the machinery
\emph{around} a model rather than its prompt. Steinberger compressed it into a slogan ---
stop prompting coding agents, design the loops that prompt them
\citep{steinberger2026loops} --- Osmani's essay named the practice \emph{loop engineering}
\citep{osmani2026loop}, and Swyx framed the general form as \emph{loopcraft}, ``the art of
stacking loops'' \citep{swyx2026loopcraft}. The stack is commonly decomposed into four
loops \citep{langchain2026loopcraft}:

\begin{description}
  \item[L1 --- Agent loop.] The model calls tools repeatedly until the task is done.
    \emph{Out of scope here}: our worker is a fixed generate--critique--refine sequence with
    no tool-calling agent.
  \item[L2 --- Verification loop.] The output is scored against a rubric and retried with
    feedback if it fails. \textbf{This is our object of study.} The \texttt{ensure\_correct}
    gene \emph{is} the verifier, the PASS gate plus \texttt{loop\_max\_retries} \emph{is}
    the stop rule, and our fix-budget sweep ($R{:}1\!\to\!3$) is a stop-rule experiment.
  \item[L3 --- Event-driven loop.] External events trigger agent runs inside a larger
    system. \emph{Out of scope}: our loop is invoked once per benchmark task, so we study no
    triggers, schedules, or cross-task memory.
  \item[L4 --- Hill-climbing loop.] Execution traces feed an analysis step that improves the
    harness configuration --- called ``arguably most important''
    \citep{langchain2026loopcraft}. \textbf{This is our intervention.} GEPA reflects on
    recorded traces and rewrites the loop's own text components; arm~D is an L4 loop pointed
    at an L2 loop.
\end{description}

So the paper's question, in this vocabulary, is: \emph{can a hill-climbing loop improve a
verification loop by rewriting its verifier?} ``Loop-level'' throughout means the verifier
and its stop rule --- not the full stack --- and we vary topology only through an optional
\texttt{plan} module and an optional repair \texttt{actuator}.

Three consequences for how we read our own results. First, this framing is what makes the
defect in Section~\ref{sec:defect} worth a paper rather than a bug report: the failure is
not GEPA-specific but a hazard for \emph{any} L4 loop, because attributing trace evidence to
the wrong L2 component optimizes the wrong thing while the scores stay plausible. Second,
Osmani's central prescription for L2 is a \emph{separated} verifier --- ``the model that
wrote the code is way too nice grading its own homework''
\citep{osmani2026loop} --- which is a qualitative claim our design measures directly: the
intrinsic critique scores $\mathrm{F}_1\!\approx\!0.11$ as a failure detector
(\S\ref{sec:analysis}), and our cross-lineage critic run tests the separation remedy.
Third, chronology: our runs were pre-registered and executed around and before these posts
appeared, so the correspondence is convergent rather than causal --- the practitioner
literature supplied the vocabulary and the stakes, not the hypotheses.

\paragraph{Reflective and evolutionary prompt optimization.}
GEPA \citep{agrawal2025gepa} reports beating GRPO by up to $20\%$ with $35\times$
fewer rollouts and exceeding MIPROv2 by over $10\%$. Crucially for us, GEPA
optimizes \emph{one or more} components and is known to plateau: practitioner
reports note GEPA returning the seed prompt essentially unchanged on
GPQA-Diamond and MATH-500. OPRO \citep{yang2023opro} and EvoPrompt use the LLM
as an optimizer; both can also \emph{degrade} performance (e.g. EvoPrompt on
GSM8K), and OPRO's authors note optimized instructions can overfit a training
split. Our null is a member of this ``plateau/degrade'' family, isolated under
controlled conditions.

\paragraph{Self-refine and the limits of self-correction.}
Self-Refine \citep{madaan2023selfrefine} popularized generate--critique--refine
with a single model, reporting $\sim\!20\%$ gains on generation-quality tasks.
Huang et al.\ \citep{huang2023cannot} show the sharp boundary: \emph{intrinsic}
self-correction (no external feedback) fails to improve and often degrades
(CommonSenseQA $75.8 \!\to\! 38.1$ after one round, recovering only to $41.8$ after a
second), and prior positive results
relied on \emph{oracle} stopping signals that ``vanish when oracle labels are
not available.'' Self-correction helps only with a genuine external verification
mechanism (a tool, a code executor, unit tests). A parallel line shows LLM
self-verifiers are unreliable detectors --- high false-positive rates
\citep{valmeekam2023selfcritique} and a self-enhancement bias, with cross-family
verification outperforming self-verification. These results predict our loop-level
null and motivate our feedback-tool and cross-verifier directions.

\paragraph{When does optimization help?}
Recent work explicitly frames the boundary question for multi-agent systems and
cost-aware settings. We adopt the same posture for the narrower loop-level vs.\
prompt-only question and contribute a pre-registered answer with controls.

\section{Method}
\label{sec:method}

\paragraph{The loop.}
The worker is a real generate--critique--refine loop. A \texttt{generate} gene
drafts an answer; a \texttt{ensure\_correct} (verify) gene drives an LLM
self-critique that returns \texttt{PASS} or a list of fixes and gates a bounded
number of refine rounds; the final answer is scored by a programmatic metric.
Critically, the critique is \emph{intrinsic}: the programmatic scorer is never
consulted inside the loop (that would be oracle contamination), so the loop
cannot game its own metric. Each eval costs $1 + 2R$ LLM calls for $R$ refine
rounds (or $2 + 2R$ with the optional plan module below).

\paragraph{Four arms.}
All arms share an identical seed genome; they differ only in what GEPA may
evolve, so a held-out gap isolates componentized optimization from a
``more prompt text'' confound.
\begin{itemize}
  \item \textbf{A} --- seed baseline, evolves nothing.
  \item \textbf{B} --- prompt-only: evolves the \texttt{generate} gene.
  \item \textbf{C} --- command control: evolves the whole genome but scores with
    a genome-blind worker; its landscape is flat by construction (a harness
    soundness check).
  \item \textbf{D} --- loop-level: evolves the whole genome
    (\texttt{generate}\,$+$\,\texttt{ensure\_correct}, and \texttt{plan} when
    enabled) with the LLM worker.
\end{itemize}
$\dd = \mathrm{mean}(D) - \mathrm{mean}(B)$ is the quantity of interest.
A and C are expected flat. \textbf{Caveat, established after the fact}
(\S\ref{sec:defect}): as originally implemented neither control could fail --- A copied its
own seed score with $\delta$ hardcoded to $0.0$, and C's command loop scored a saturated
$1.0$, so it proposed nothing. The corrected harness re-measures A and gives C a
genome-blind constant draft scored by the real metric, which is flat without being
saturated.

\paragraph{Optimizer.}
We drive the official GEPA \citep{agrawal2025gepa} via a loop-genome adapter:
\texttt{gepa.optimize(seed\_candidate, trainset, adapter, reflection\_lm,
max\_metric\_calls)} with a \texttt{MaxCandidateProposalsStopper} to make the
reflection-LM spend a hard cap. GEPA's reflection channel (offline, between
training iterations) sees per-constraint feedback prose --- this is the standard
optimizer signal and is identical across B and D; it is distinct from the
in-loop critique gate.

\paragraph{Two levers.}
(i) \emph{Feedback richness}: the critique step is optionally given
\texttt{text\_stats(draft)}, a deterministic one-line readout of measured
properties of the draft (word/character/sentence/bullet counts, longest-word
length, last words). It takes only the text and never sees the task, so it is a
\emph{measurement}, the tool/executor signal \citep{huang2023cannot} found
effective --- not the oracle. (ii) \emph{Complexity}: an optional third
\texttt{plan} gene turns the loop into plan--generate--critique--refine, giving
D a third interacting module to co-tune (the DSPy lever); B still evolves only
\texttt{generate}.

\paragraph{Oracle taxonomy (a required distinction).}
Structural non-oracle-ness (the tool takes no task argument) is necessary but
\textbf{not sufficient}. A statistic is \emph{informationally} the oracle when
its verifier has no per-task parameter, so the raw measurement alone resolves
the verdict. A pre-spend design review caught exactly this: reporting
\texttt{all\_lowercase} or a comma count would hand back the verdicts of the
parameter-free \texttt{all\_lowercase} and \texttt{no\_comma} constraints. Both
statistics were removed before any paid run. Every reported statistic requires a
task-supplied target (an exact/min count, a required bullet count, a max length)
to become pass/fail, so the model must still know each target and reason --- the
legitimate tool regime.

\paragraph{Reproducibility and budget.}
Every LLM call is recorded to a content-addressed cassette (\texttt{sha256} of
name$+$args), so runs replay offline at zero cost and partial spend is always
persisted. Reproduction is a first-class, strict operation: re-running any study
with \texttt{-{}-replay} re-executes the entire pipeline (loaders, GEPA search,
scoring) against the cached calls with no API key, and any cassette miss raises.
GEPA is deterministic under replay, so replay is bit-for-bit:
\texttt{scripts/verify\_reproducibility.py} registers eleven committed runs and confirmed
that every per-arm held-out score matched its released JSON exactly. \textbf{Those eleven
runs are now labelled superseded} (\S\ref{sec:defect}): the corrected harness changes every
reflection prompt and therefore every cassette key, so it cannot replay them, and the script
reports them as skipped rather than verified. Output tokens are hard-capped via enforced \texttt{max\_tokens};
a pre-flight guard refuses any run whose (conservative) estimate exceeds
\texttt{-{}-max-usd}, whose default is \$5; three runs raised it explicitly (to \$6.18,
\$5.89 and \$11.76), so \$5 is a default rather than a study-wide ceiling. Every experiment was pre-registered --- hypothesis, frozen
configuration, expected seed band, and verdict rule committed to version control
\emph{before} spending.

\paragraph{Verdict rule (pre-registered).}
\textbf{SUPPORTED} iff $\dd > \sigma_B + \sigma_D$ and positive;
\textbf{NOT SUPPORTED} iff $\mathrm{mean}(D) \le \mathrm{mean}(B)$;
\textbf{INCONCLUSIVE} if positive-but-within-noise, fewer than two seeds, or a
ceiling/floor flag fires (seed held-out $\ge 0.9$ or $\le 0.1$). This is a
deliberately conservative bar (the sum of two standard deviations), erring
toward under-claiming.

\section{Experimental Setup}
\label{sec:setup}

\paragraph{Models.} Worker \texttt{gpt-4.1-mini} (output cap 320 tokens);
reflection \texttt{gpt-5.6-terra} (output cap 1500).

\paragraph{Tasks.} A composable IFEval-style \citep{zhou2023ifeval}
multi-constraint generator with nine independent, parameter-carrying verifiers
(e.g.\ \texttt{exact\_word\_count}, \texttt{forbidden\_word},
\texttt{word\_max\_length}, \texttt{num\_bullets}). Joint satisfiability of every
generated task is guaranteed by a witness the generator validates to score
$1.0$ (verified exhaustively: $3{,}189$ easy-tier and $19{,}632$ hard-tier
combinations). The \emph{hard} tier mandates \texttt{exact\_word\_count} (a
genuine LLM skill gap) plus at least one further ``brutal'' constraint, which
structurally caps the seed below saturation.

\paragraph{Metric.} Mean fraction of constraints satisfied on a disjoint
held-out split (\texttt{mc\_score} $\in [0,1]$), scored programmatically (no LLM
judge). We report a relative result (D vs.\ B under matched budget), not a
leaderboard number.

\section{Results}
\label{sec:results}

Table~\ref{tab:results} collects all ten pre-registered runs. Read them under the audit in \S\ref{sec:defect}: the $\dd(D-B)$ column does not measure what its header claims, because the loop-level arm's extra gene was never changed.

\begin{table}[t]
\centering
\small
\begin{tabular}{clccccl}
\toprule
Run & Regime & Seeds & Seed & B (prompt-only) & D (loop-level) & \dd $=D-B$ \\
\midrule
1 & single-constraint            & 2 & 0.47 & 0.59 & 0.50 & $-0.09$ \\
2 & easy multi ($k{=}3{-}4$)     & 5 & \textbf{0.97} & 0.97 & 0.97 & $0.00$ \\
3 & hard multi ($k{=}4{-}5$)     & 5 & 0.765 & 0.770 & 0.770 & $\mathbf{0.00}$ \\
4 & hard multi $+$ feedback tool & 3 & 0.837 & 0.798 & 0.798 & $\mathbf{0.00}$ \\
5 & hard multi $+$ plan module   & 3 & 0.757 & 0.757 & 0.757 & $\mathbf{0.00}$ \\
6 & hard multi $+$ cross-lineage critic & 3 & 0.825 & 0.817 & 0.825 & $+0.01$ (INCONCL.) \\
7 & hard multi $+$ oracle critique ($R{=}1$) & 3 & \textbf{0.748} & 0.74 & 0.74 & $\mathbf{0.00}$ \\
8 & hard multi $+$ oracle critique ($R{=}3$) & 3 & \textbf{0.813} & 0.80 & 0.80 & $\mathbf{0.00}$ \\
9 & hard multi $+$ oracle ($R{=}3$), feedback \emph{lo} (verdict) & 3 & \textbf{0.810} & 0.80 & 0.80 & $\mathbf{0.00}$ \\
10 & hard multi $+$ oracle ($R{=}3$), feedback \emph{hi} (edits) & 3 & \textbf{0.873} & 0.86 & 0.86 & $\mathbf{0.00}$ \\
\bottomrule
\end{tabular}
\caption{All ten pre-registered runs. Runs 1--2 are uninformative by design
(the verify gene is inert on single-constraint tasks; the seed saturates at
$0.97$ on easy multi). Runs 3--10 are fair tests with genuine headroom and no
ceiling/floor flag. In every one $\dd \le 0.01$; the sole non-negative gap
(run 6, a different-model critic) is $+0.01$, within seed noise
(\textbf{INCONCLUSIVE}), and every other fair test is $\dd = 0.00$
(\textbf{NOT SUPPORTED}). Runs 7--10 are labeled ceiling/mechanism controls
(\S\ref{sec:analysis}): the seed is the readout, not \dd. Runs 8--10 share the
oracle $R{=}3$ regime and vary only the repair-feedback actionability (mid $=$
run 8): terse verdict (lo, $0.810$) $\approx$ which-constraint (mid, $0.813$)
$\ll$ concrete measured edits (hi, $0.873$).}
\label{tab:results}
\end{table}

\paragraph{A different-model critic nudges, but not significantly (run 6).} Giving
the critique step a different-lineage model (\texttt{gpt-5.6-luna} verifying a
\texttt{gpt-4.1-mini} worker) --- motivated by the finding that same-family judges
carry a self-preference bias --- is the first run where D does not lose: $\dd =
+0.01$. But at $n=3$ this is within the pooled seed noise, so the pre-registered
verdict is INCONCLUSIVE. The critic also raised the seed baseline ($0.765 \to
0.825$). This is directionally consistent with the reliable-critic hypothesis and
motivates a properly-powered, \emph{true} cross-family (cross-provider) follow-up;
we do not claim an effect. \textbf{Caveat:} \texttt{gpt-5.6} and \texttt{gpt-4.1}
share a vendor and training lineage, so this is cross-\emph{lineage}, not a true
cross-\emph{family} test --- residual self-preference bias may remain.

\paragraph{A perfect detector does not lift the ceiling (run 7).} The oracle-critique
control --- where the in-loop gate is the true metric --- was pre-registered to test
whether the loop \emph{can} convert a perfect detector into higher performance
(expected seed $\ge 0.85$). It cannot: the oracle seed is $0.748$, level with the
intrinsic baseline ($0.765$) and \emph{below} the tool baseline ($0.837$). Detection
quality is therefore \textbf{not} the binding constraint. Even told exactly which
constraints fail, \texttt{gpt-4.1-mini} cannot reliably fix hard constraints (e.g.
an exact word count) within the bounded refine budget. This is the crux
(\S\ref{sec:analysis}).

\paragraph{More fix attempts help, but saturate (run 8).} Holding the detector perfect
(oracle) and raising the refine budget $1 \to 3$ lifts the seed $0.748 \to 0.813$: the
\emph{fix} step is a genuine lever. But it still falls short of the pre-registered $0.85$
bar and below the tool baseline ($0.837$) --- a good tool at one retry roughly equals the
oracle at three. The residual is the generator's raw capability: perfect, repeated feedback
cannot push \texttt{gpt-4.1-mini} past $\sim\!0.81$ on hard constraints. The gains a
self-refine loop can extract (better detection $+$ more attempts) are real but bounded, and
prompt-only optimization already captures them. This is consistent with contemporaneous
code-generation findings that the feedback signal and the number of repair attempts --- not
the loop's structure --- drive self-repair gains \citep{feedbackoverform2026,howmanytries2026}.

\paragraph{Feedback \emph{actionability} is a real fix-side lever --- but it is feedback
shape, not GEPA (runs 9--10).} Runs 7--8 leave the residual fix-step gap ambiguous: is
the worker unable to \emph{act} on good feedback (generator-limited), or is the feedback
itself insufficiently actionable (feedback-limited)? We separate them by holding detection
fixed at the perfect oracle ($R{=}3$) and varying only the repair text handed to the refine
step across three pre-registered levels: \emph{lo} (a bare ``revise'' verdict), \emph{mid}
(which constraints failed, with their targets --- identical to run 8), and \emph{hi} (mid
\emph{plus} a concrete \emph{measured} edit per failed constraint, e.g.\ ``the draft has 2
words but needs exactly 20: add 18''). The primary contrast is decisive: $\dd(\text{hi}-
\text{lo}) = 0.873 - 0.810 = +0.063$ on the seed held-out (arm A, no GEPA), clearing the
pre-registered $+0.05$ threshold --- the fix step is partly \textbf{feedback-limited}.
All three seeds are positive ($+0.090, +0.025, +0.075$) and an exact-resample $95\%$
interval $[+0.025, +0.090]$ excludes zero, so the effect survives at $n{=}3$ --- though a
sign test at $n{=}3$ cannot fall below $p=0.25$, so this rests on the magnitudes rather than
on replication across seeds, and the whole comparison is a ceiling result measured with the
true metric gating at inference (\S\ref{sec:defect}). Two
sharp qualifications follow. (i) The gain is entirely in the \emph{hi} step: lo $\approx$
mid ($0.810$ vs.\ $0.813$), so naming \emph{which} constraint failed adds nothing over a
bare verdict; only the concrete, measured edit moves the needle. (ii) The lever is the
feedback's \emph{shape}, a hand-built measurement function --- \emph{not} loop-level GEPA:
$\dd(D-B)=0.00$ at every level, and optimization again slightly \emph{degrades} both arms
(hi seed $0.873 \to$ opt $0.86$). Finally, \emph{hi} ($0.873$) is the first oracle
configuration to \emph{exceed} the tool baseline ($0.837$), resolving the run-4 tool
$>$ oracle inversion (\S\ref{sec:analysis}) as an actionability effect: the oracle
underperformed the tool at $R{=}1$ because its verdicts were terse, not because a perfect
detector hurts. Yet even maximal, exact repair guidance saturates at $0.873$ (still short of
perfect, no ceiling flag) --- so a generator-limited residual remains. The fix step is thus
\emph{both} feedback-limited (up to $\sim\!+0.06$ from actionability) and, beyond that,
generator-limited.

\paragraph{The fair tests agree.} Run 3 is the decisive baseline: seed held-out
$0.765$ sits inside the pre-registered band, so the loop had real headroom and
the verify gene had multiple constraints to decompose --- yet $\dd = 0.00$. The
tie is not a flat line but a wash: per-seed, D and B trade wins
(e.g.\ seed 1 $D\!+\!0.05$; seed 4 $B\!+\!0.055$) and net to zero. Control C is
dead-flat on every seed --- which we read at the time as confirming the harness, and which
\S\ref{sec:defect} shows it cannot, because that control was saturated and therefore
incapable of moving. An exact-resample (all $n^n$ resamples) $95\%$ interval on the paired
per-seed $\dd$ is $[-0.033, +0.033]$ for run 3 ($n{=}5$) and \emph{exactly}
$[0.000, 0.000]$ for runs 4 and 10. \textbf{We previously called this ``tight, not merely
underpowered''; that reading is withdrawn} (\S\ref{sec:defect}). The $[0,0]$ runs are ones
where B and D produced byte-identical optimized genomes on every seed, so both arms issued
the same calls and the cassette served the same completions: the zero is an identity, and the
interval is evidence that the treatment was not applied rather than evidence that it did
nothing. Of the four fair null runs, $3/5$, $3/3$, $2/2$ and $2/3$ seeds are such ties.

\paragraph{Feedback helps the baseline, not the gap (run 4).} Adding the
measurement tool raised the seed loop from $0.765$ (run 3) to $0.837$ --- direct
confirmation of the external-feedback prediction \citep{huang2023cannot}. But
optimization then \emph{degraded} both arms (mean $-0.04$; $-0.10$ on one seed),
and D tracked B exactly on all three seeds. The tool helps B and D equally, so it
is not evidence for loop-level optimization.

\paragraph{A third module gives no traction (run 5).} With a \texttt{plan}
module (D co-evolving three genes, B still one), optimization moved nothing on
any arm and $\dd = 0.00$. More optimizable structure did not translate into a
loop-level advantage at this budget.

\subsection{Cross-domain corroboration: BigCodeBench-Hard}
\label{sec:bcb}

To test whether the instruction-following null generalizes to a domain where the
verify step catches real \emph{semantic} errors and runs as an executable tool ---
the regime \citet{shinn2023reflexion} find self-correction works --- we replicated
the four-arm design on \textbf{BigCodeBench-Hard} \citep{zhuo2024bigcodebench}, hard
hard enough that strong models score far below their full-benchmark pass@1 on it (the
Hard subset is calibrated to stay challenging) \citep{zhuo2024bigcodebench}. (An earlier attempt on MBPP+
\citep{liu2023evalplus} saturated its training split and is reported in
Appendix~\ref{app:mbpp-failed}; that failure motivated the move to a frontier-hard
benchmark.) Each task's \texttt{unittest} methods are partitioned into a visible half
(the in-loop, executable signal) and a disjoint hidden half (the held-out scorer),
reconstructing the anti-oracle split; solutions are graded in a pinned, offline Docker
sandbox (\texttt{-{}-network none}, memory/PID caps, single-thread BLAS).\footnote{Executing
untrusted \emph{model} code at scale surfaces a subtlety absent when only reference solutions
run: a draft can exhaust the container's memory cap. Because we run one capped container at a
time on a host with ample free RAM, such an out-of-memory kill is the \emph{solution's} failure
(scored $0.0$), not a host fault --- distinct from a genuine infrastructure fault (daemon down,
image/mount error), which aborts loudly rather than silently recording a $0.0$.} We ran arms
A--D on $2$ seeds, $8$ train / $10$ held-out tasks drawn from a sandbox-verified
$80$-task pool, worker \texttt{gpt-4.1-mini}, reflection \texttt{gpt-5.6-terra}, at
$\mathrm{mmc}=30$.

The seed baseline sits at $0.40$ (seed 0) and $0.50$ (seed 1) --- far from both floor
and ceiling, so the loop has genuine headroom, unlike MBPP+. And unlike MBPP+, GEPA
\emph{ran}: on seed 1 it accepted improving candidates for \emph{both} the prompt-only
arm B (evolving \texttt{generate}) and the loop-level arm D (evolving \texttt{generate}
\emph{and} \texttt{ensure\_correct}), each lifting the $8$-task training aggregate
$0.375 \to 0.625$; on seed 0 neither arm accepted a candidate within budget, and the
genome-blind control C is flat by construction. Yet on the held-out split \emph{every
arm scored exactly its seed baseline} ($0.40$, $0.50$): the accepted training
improvements changed which held-out tasks \texttt{gpt-4.1-mini} solves not at all. The
headline is therefore a clean $\dd(D-B)=0.00$ (exact-resample $95\%$ CI
$[0.000, 0.000]$) --- and, more pointedly, even when loop-level optimization has its
extra lever and \emph{pulls} it (an accepted \texttt{ensure\_correct} rewrite that
improves training), it yields no held-out gain over prompt-only. The
instruction-following null recurs in a semantically-verified, executable-feedback
domain: train-split improvement does not transfer, and the loop-level lever is inert
on held-out.

\paragraph{Actionable feedback lifts the code baseline too (cross-domain D2$'$).} The one
intervention that helped on instruction-following --- concrete, measurement-shaped repair
feedback (runs 9--10) --- transfers to code, and \emph{more strongly}. Holding the arm at A
(no GEPA, so the lever is feedback \emph{shape}, not optimization) and varying only the in-loop
visible-test feedback between a bare verdict (\emph{lo}) and the verdict plus the source of the
visible tests the solution must satisfy (\emph{hi}), the held-out (hidden-test) score rises by
a mean $+0.20$ across six seeds (per-seed $+0.40, 0.00, +0.20, +0.20, +0.20, +0.20$). An
exact-resample $95\%$ interval is $[+0.100, +0.300]$, \emph{excluding zero} --- a
$3\times$-larger effect than the instruction-following actionability lift ($+0.063$), plausibly
because a code assertion is a sharper repair target than a word-count. Two qualifications hold
as on instruction-following: the gain accrues to the \emph{baseline}, not to $\dd(D-B)$ (which
stays $0.00$), so it is feedback \emph{shape}, not loop-level optimization --- GEPA did not need
to discover it; and the held-out scores carry some execution non-determinism (the effect size
dwarfs it, but we do not over-read the exact value). This reinforces the paper's conclusion:
the productive lever is the fixer and the actionability of its feedback, not loop-level
optimization of the critique.

\paragraph{Can GEPA \emph{discover} the actionable feedback? The actuator experiment.}
The $+0.20$ lift establishes that the fixer's feedback \emph{matters} --- but that feedback
was hand-designed. The sharpest test of our thesis is whether GEPA can \emph{discover} the
gain by optimization when pointed squarely at it. We add a fourth evolvable gene, an
\texttt{actuator}: an LLM step that, on a failed critique, turns the visible-test signal and
the visible-test source (the same \emph{hi} material) into the repair guidance fed to the
refine step. Arm~B still evolves only \texttt{generate}; arm~D now co-evolves \texttt{generate}
\emph{and} \texttt{actuator} with the detector (\texttt{ensure\_correct}) held \emph{fixed}, so
any $D-B$ gap isolates optimization of the fixer's feedback --- the one component we know can
help. A first run ($2$ seeds, $8$/$10$) gave a flat $\dd(D-B)=0.00$: GEPA proposed sensible
actuator rewrites (e.g.\ ``produce concrete, draft-specific repair instructions; do not write
vague advice such as `fix the syntax' ''), but accepted none that beat even the training
subsample, and the single accepted \texttt{generate} rewrite did not transfer.

To rule out that this null was a small-budget artifact, we ran a \emph{powered} replication:
$3$ seeds, $20$ train / $20$ held-out, $\mathrm{mmc}=50$ --- $2.5\times$ the train size,
$2\times$ the held-out resolution, $1.7\times$ the search. The powered result does not merely
reproduce the null; it sharpens \emph{against} loop-level optimization. Prompt-only arm~B now
\emph{transfers} to held-out ($+0.05, +0.15, 0.00$; mean $+0.067$) --- with a larger training
split, evolving \texttt{generate} alone generalizes. But the loop-level arm~D does \emph{worse}:
$\dd(D-B) = -0.10$ (per-seed $-0.15, -0.15, 0.00$; exact-resample $95\%$ CI
$[-0.150, +0.000]$, which \emph{includes} zero; sign test $p = 0.50$ on the two non-tied
seeds), and on seed~0 arm~D's optimization actively \emph{degrades} held-out
($0.55 \to 0.45$). \textbf{Two corrections apply to this paragraph as first written}
(\S\ref{sec:defect}). We reported the interval as $[-0.150,-0.050]$ and ``excluding zero'';
that was a percentile-index error, and corrected it includes zero. And we explained the deficit
as textbook over-fitting --- arm~D's larger search space fitting the $20$ training tasks harder
--- which the harness contradicts: under GEPA's default one-component-per-proposal selector,
arm~D searched the shared \texttt{generate} gene roughly \emph{half} as hard as arm~B, and the
train aggregate that would have distinguished the two explanations was computed and discarded.
What remains sayable is the ordering: across all three code configurations $D \le B$ on
held-out, never $D > B$; and on $20$ binary-scored held-out tasks $-0.10$ is two task flips,
against a seed spread of four to five flips across configurations. The honest reading is that
raising train size, held-out resolution, seeds and search depth surfaced no hidden positive
\emph{for an arm whose extra gene the harness was not equipped to optimize}.

\paragraph{Reproducibility scope.} The main four-arm result reproduced \emph{identically}
across two independent paid runs. Per-execution grading is deterministic, and we release the
cassette and scores. We do \emph{not} claim byte-for-byte replay of the full LLM-call
trace here, however: model solutions may use unseeded randomness, so re-executing them
in the sandbox can yield run-to-run-varying test output --- an intrinsic property of
executable-code benchmarks, not of our harness. (The instruction-following runs above,
which never execute model code, do replay bit-for-bit.) We normalize the string-level
non-determinism we can --- object ids, addresses, and temp paths in failure messages ---
so that identical executions render identical feedback. The actuator runs (both the $2$-seed
and the powered $3$-seed replication) share this scope: cassettes and scores released,
per-execution grading deterministic, full-trace replay not claimed.

\section{Analysis: Why the Loop-Level Lever Is Inert}
\label{sec:analysis}

We advance three converging explanations.

\paragraph{(1) The critique is an unreliable detector --- measured.} The
mechanism by which a better verify gene \emph{could} help is by catching failures
the generator cannot self-see and steering refinement. We measure the in-loop
critique directly as a binary classifier of ``the draft satisfies all
constraints,'' computed offline from the recorded cassettes (zero cost) by
pairing each critique's \texttt{PASS}/\texttt{FAIL} verdict with the ground-truth
\texttt{mc\_pass} of the exact draft it judged, on the seed loop over the
held-out split (Table~\ref{tab:detector}). The intrinsic critique is a
\emph{near-random} detector: F$_1 = 0.11$ on run~3 and effectively $0$ on run~5
(it emits \texttt{PASS} on drafts that fail, precision $0.08$). A detector this
poor offers no reliable signal to sharpen, so evolving its rubric cannot
systematically improve outcomes --- the direct mechanism of the loop-level null.
Strikingly, the \emph{same} critique fed the measurement tool (run~4) becomes a
reliable detector (precision $1.00$, F$_1 = 0.83$) --- which is precisely why the
tool lifts the baseline (Section~\ref{sec:results}); but that reliable detector
sits in \emph{both} the B and D loops, so it does not open a loop-level gap. This
matches the documented unreliability of LLM self-verification --- high
false-positive rates, self-enhancement bias, cross-family verifiers beating
self-verification \citep{valmeekam2023selfcritique,huang2023cannot}.

\begin{table}[h]
\centering
\small
\begin{tabular}{lccccc}
\toprule
Critique regime & PASS-precision & recall & F$_1$ & accuracy & $n$ \\
\midrule
run 3 (intrinsic, no tool)   & 0.08 & 0.20 & 0.11 & 0.68 & 50 \\
run 5 (intrinsic $+$ plan)   & 0.00 & 0.00 & $\approx 0$ & 0.70 & 30 \\
run 4 (tool-fed critique)    & \textbf{1.00} & 0.71 & \textbf{0.83} & 0.93 & 30 \\
\bottomrule
\end{tabular}
\caption{The in-loop critique as a failure detector (seed loop, held-out split,
from cassettes). Intrinsic self-critique is near-random; a non-oracle measurement
tool makes it reliable. The tool-fed reliability helps B and D equally.}
\label{tab:detector}
\end{table}

\paragraph{(1b) But detection quality is necessary, not sufficient --- the
non-monotonic ceiling.} If a poor detector were the \emph{only} obstacle, a better
one should monotonically raise the loop's held-out score. It does not
(Table~\ref{tab:curve}). The intrinsic critique (F$_1\!=\!0.11$) yields seed
$0.765$; the tool-fed critique (F$_1\!=\!0.83$) yields $0.837$; but the
\emph{oracle} critique (F$_1\!=\!1.0$ by construction) yields only $0.748$ --- the
best detector gives the \emph{worst} of the three. Detection quality is therefore
necessary but not sufficient. Two non-exclusive mechanisms explain the inversion,
both pointing past detection: (i) the \emph{fix} step is the true bottleneck ---
even told exactly which constraints fail, the worker cannot reliably repair a hard
constraint (an exact word count) in one refine round; and (ii) \emph{feedback
actionability} matters more than feedback correctness --- the tool-fed LLM critique
emits natural, instructive fix prose, whereas the oracle emits terse per-constraint
verdicts, so the tool-fed refine is better guided despite a less-than-perfect
detector. Runs 9--10 test (ii) directly (\S\ref{sec:results}): holding detection at
the perfect oracle and making the repair feedback \emph{maximally} actionable
(concrete measured edits) lifts the seed to $0.873$, \emph{above} the tool baseline
($0.837$) and $+0.063$ over a terse verdict --- so the tool$>$oracle inversion was an
actionability artefact, not evidence that a perfect detector hurts. Mechanism (i)
also survives: even the maximally-actionable oracle saturates at $0.873$, so a
generator-limited residual remains. And the pre-registered ceiling verdict is
decisive either way --- no oracle configuration reaches a loop-level advantage
($\dd(D-B)=0.00$ throughout).

\begin{table}[h]
\centering
\small
\begin{tabular}{llccc}
\toprule
Critique detector & refine budget & detector F$_1$ & seed held-out & run \\
\midrule
intrinsic self-critique & 1 & 0.11 & 0.765 & 3 \\
non-oracle measurement tool & 1 & 0.83 & 0.837 & 4 \\
oracle, which-constraint (mid) & 1 & 1.00 & 0.748 & 7 \\
oracle, which-constraint (mid) & 3 & 1.00 & 0.813 & 8 \\
oracle, bare ``revise'' (lo) & 3 & 1.00 & 0.810 & 9 \\
oracle, concrete measured edits (hi) & 3 & 1.00 & \textbf{0.873} & 10 \\
\bottomrule
\end{tabular}
\caption{Detector quality vs.\ loop performance is \textbf{non-monotonic}: the
perfect (oracle) detector with which-constraint feedback (mid; runs 7--8) does not
maximize held-out score, falling below a good tool at one retry (run 4). Two fix-side levers close the
gap: raising the refine budget ($0.748\!\to\!0.813$, run 8) and --- decisively ---
making the feedback \emph{actionable} (concrete measured edits, run 10: $0.873$, now
\emph{above} the tool). Detector F$_1$ is fixed at $1.0$ across runs 7--10, so the
$0.810\!\to\!0.873$ lift is pure feedback \emph{actionability}, not detection. A
generator-limited residual remains (hi still $<1$).}
\label{tab:curve}
\end{table}

\paragraph{(1c) The residual is one constraint, not a diffuse ceiling.} We decompose
the fix-step residual by constraint: replaying each run's cassette (offline, \$0;
\texttt{scripts/analyze\_residual\_by\_constraint.py}), we take the loop's \emph{final}
draft for every held-out task and record the per-constraint satisfaction rate
(Table~\ref{tab:residual}). The result is stark. Under the oracle detector with
maximally-actionable feedback (run 10), every ``choice'' constraint --- \texttt{no\_comma},
\texttt{all\_lowercase}, \texttt{num\_bullets}, \texttt{include\_keyword},
\texttt{end\_with}, \texttt{forbidden\_word} --- is satisfied at rate $1.0$, and
\texttt{word\_max\_length} at $0.90$. The residual is almost entirely
\texttt{exact\_word\_count} (with a small remainder in \texttt{word\_max\_length}):
even told exactly ``the draft has 2 words but needs exactly 20: add 18,''
\texttt{gpt-4.1-mini} hits the exact count only $\sim\!0.50$ of the time (up from
$0.16$ intrinsic, $0.33$ with which-constraint feedback). This is the known
inability of LLMs to control exact token counts during generation. So the
generator-limited residual is not a broad capability wall --- it is a single mechanical
skill, of a kind a deterministic actuator (not prompt/critique optimization) would
address; the loop-level lever is inert regardless.

\begin{table}[h]
\centering
\small
\begin{tabular}{lccc}
\toprule
Constraint (final-draft satisfaction) & run 3 (intrinsic $R1$) & run 8 (oracle mid $R3$) & run 10 (oracle hi $R3$) \\
\midrule
\texttt{exact\_word\_count}  & 0.16 & 0.33 & \textbf{0.50} \\
\texttt{word\_max\_length}   & 0.85 & 0.80 & 0.90 \\
\texttt{forbidden\_word}     & 0.92 & 0.96 & 1.00 \\
\texttt{num\_bullets}        & 1.00 & 1.00 & 1.00 \\
\texttt{end\_with}           & 0.95 & 1.00 & 1.00 \\
\texttt{include\_keyword}    & 1.00 & 1.00 & 1.00 \\
\texttt{all\_lowercase}      & 0.95 & 1.00 & 1.00 \\
\texttt{no\_comma}           & 0.95 & 1.00 & 1.00 \\
\bottomrule
\end{tabular}
\caption{Per-constraint satisfaction of the loop's final draft (held-out, from
cassettes). The fix-step residual is concentrated almost entirely in
\texttt{exact\_word\_count}; actionable feedback lifts it ($0.16\!\to\!0.33\!\to\!0.50$)
but never solves it, while every choice constraint reaches $1.0$. The generator-limited
residual is one mechanical skill (exact counting), not a diffuse ceiling.}
\label{tab:residual}
\end{table}

\paragraph{(2) Intrinsic self-correction has a low ceiling.} Consistent with
\citet{huang2023cannot}, the run-4 result shows the loop improving when given an
external tool (seed lift) but the \emph{reflective optimization on top} adding
noise, not signal --- the ``self-correction can make it worse'' regime.

\paragraph{(3) Optimization overfits the training split.} Run 4's optimized
held-out scores fall \emph{below} the seed, echoing OPRO's observation that
optimized instructions can overfit a training set \citep{yang2023opro}. With a
small budget and small splits, GEPA's proposals can win on the training
subsample and lose on held-out; this hurts B and D alike.

\section{A Defect in Our Own Harness, and What It Invalidates}
\label{sec:defect}

An adversarial audit of the released artifact, run after the experiments were complete,
established that the loop-level arm had not received its intervention. We state the evidence,
then the retractions, then the fix. This section supersedes any reading of
Section~\ref{sec:results} as a test of loop-level optimization.

\paragraph{Evidence 1: the critique gene never changed.}
Replaying run~3 --- our flagship $n{=}5$ multi-hard run --- offline from its cassette and
dumping every winning genome gives: \texttt{generate} changed in $5/15$ winners,
\texttt{ensure\_correct} in $\mathbf{0/15}$. GEPA \emph{did} propose critique rewrites and
they were rejected; more revealingly, the proposals were \emph{generation} instructions
written into a critique slot, e.g. ``\emph{For every user request, produce an answer that
satisfies the requested content task and every stated formatting or linguistic constraint
exactly}''. Whatever arm~D measured, it was not a loop with an optimized verifier.

\paragraph{Evidence 2: the reflector could not see the critique.}
The cause is a single line: our reflective dataset returned the \emph{same} records for every
component, and those records' ``Generated Outputs'' field was defined as the last
\emph{worker} step --- deliberately excluding verifier steps. So when the optimizer reflected
on the verify gene it was shown the task and the \emph{generator's} answer, with no
information about what the critique said, whether it fired, or whether its verdict was right.
A detector cannot be improved from evidence that never mentions its verdicts, and the
observed behaviour (proposing generation prose into a critique slot) is exactly what a
reflector shown only generator I/O should be expected to produce.

\paragraph{Evidence 3: the two-gene arm was budget-starved per gene.}
We passed no component selector to GEPA, which defaults to selecting one component per
proposal and cycling. Under a matched \emph{total} proposal cap, a two-gene arm therefore
searches each gene about half as hard as a one-gene arm. Measured proposal targets on run~3
confirm it: arm~B spent up to four proposals on \texttt{generate}, arm~D about two, spending
the rest on the gene that never got accepted. This also reverses the mechanism we had
asserted --- we wrote that arm~D ``fits the training tasks harder'' and so transfers worse,
whereas under this selector arm~D fits the shared gene \emph{less} hard.

\paragraph{Evidence 3b: half the headline effect is visibly a budget artifact.}
The powered actuator run's per-seed held-out scores, from the released JSON and requiring no
replay, decompose $\dd(D-B)=-0.10$ into two unlike halves:

\begin{center}\small
\begin{tabular}{cccc}
\toprule
seed & B (seed $\to$ opt) & D (seed $\to$ opt) & $\dd(D-B)$ \\
\midrule
0 & $0.55 \to 0.60$ & $0.55 \to \mathbf{0.45}$ & $-0.15$ \\
1 & $0.45 \to \mathbf{0.60}$ & $0.45 \to 0.45$ \emph{(unmoved)} & $-0.15$ \\
2 & $0.65 \to 0.65$ & $0.65 \to 0.65$ & $0.00$ \\
\bottomrule
\end{tabular}
\end{center}

A score that moves proves the winning genome differed from the seed, so on seed~0 arm~D
\emph{did} ship a change --- one that transferred badly. The audit's ``treatment never
applied'' verdict therefore does not extend uniformly to this experiment, which evolves
\texttt{generate} and \texttt{actuator} with the detector held fixed; which of those two
genes changed on seed~0 we cannot recover, because bcb replay depends on container test
outcomes that are not reproducible across hosts (one visible test times out under load on
our machine, changing the critique prompt and hence the cassette key).

Seed~1 is the informative one. There arm~B accepted a change that transferred ($+0.15$)
while arm~D's held-out score never moved from its seed at all. That is the round-robin
signature: under a matched proposal cap, B spent its budget on \texttt{generate} and
accepted a winner, while D divided the same budget across two genes and accepted nothing
that transferred. Half of the headline $-0.10$ is thus not a statement about optimizing the
fixer; it is a statement about how much search each arm was allowed to spend on the gene
they share. Note the asymmetry of what scores can prove: an unmoved score does \emph{not}
establish an unchanged genome (a changed genome can score identically), so seeds~1 and~2
remain formally inconclusive on that point.

\paragraph{Evidence 4: a percentile-index error.}
Our exact-resample interval took the upper bound at index
$\lfloor(1-\alpha/2)N\rfloor - 1$. The resample distribution is discrete with $N=n^n$
equally-likely atoms, so at $n{=}3$ each atom carries $1/27 = 3.70\% > \alpha/2 = 2.5\%$ and
the largest atom cannot be dropped. Correctly indexed at $\lceil(1-\alpha/2)N\rceil - 1$:

\begin{table}[h]
\centering
\small
\begin{tabular}{lccc}
\toprule
Result & per-seed & as published & corrected \\
\midrule
Actuator, powered ($n{=}3$) & $-0.15, -0.15, 0.00$ & $[-0.150,-0.050]$ ``excludes zero''
  & $\mathbf{[-0.150,+0.000]}$ \emph{includes zero} \\
Actionability, hi$-$lo ($n{=}3$) & $+0.090, +0.025, +0.075$ & $[+0.025,+0.085]$
  & $[+0.025,+0.090]$ (still excludes zero) \\
Null, run~3 ($n{=}5$) & $0,+0.055,0,0,-0.055$ & $[-0.033,+0.033]$ & unchanged \\
\bottomrule
\end{tabular}
\caption{The upper bound was too low by one atom. Only $n{=}3$ intervals are affected, and
the direction of the error flatters negative effects: it turns an interval that brackets
zero into one that appears to exclude it.}
\label{tab:ciexcl}
\end{table}

\paragraph{Evidence 5: the flatness controls could not fail.}
Arm~A copied its own seed score into its optimized slot with $\delta$ hardcoded to $0.0$, so
it had no second measurement. Arm~C's genome-blind command loop always exits $0$, which our
scalarization maps to exactly $1.0$ --- the optimizer's perfect score --- so it proposed
nothing and never scored a mutated genome. Our claim that ``a non-flat A or C would
invalidate the run'' was therefore unfalsifiable by construction.

\paragraph{What we retract.}
The following claims are withdrawn: that there is \emph{no measurable advantage for
loop-level optimization} (the loop-level treatment was not applied); that the null is
\emph{tight rather than underpowered} on the strength of $[0,0]$ intervals (an interval of
$[0,0]$ arises when B and D produce byte-identical genomes, i.e.\ when the arms issued the
same calls and the cassette served the same completions --- the sharpest available evidence
that the experiment did not run, not evidence for its conclusion); that the extra evolvable
gene \emph{over-fits the train split} (the train aggregate was computed and discarded, so
this was untestable from the release, and the selector makes the opposite prediction); that
$\dd(D-B)=-0.10$ \emph{excludes zero}; and that this \emph{rules out under-powering}. We also
correct two factual slips: the reproducibility script registers eleven runs, not ten, and the
\$5 ceiling is a CLI default that three runs overrode (up to \$11.76).

\paragraph{What survives.}
Everything that does not depend on arm~D having been optimized. The intrinsic critique's
detector quality ($\mathrm{F}_1\!=\!0.11$, and $0.83$ with the measurement tool), the
per-constraint residual decomposition, the fix-budget effect, and the feedback-actionability
effects ($+0.063$ text, $+0.20$ code) are all measured on arm~A or on paired conditions that
never involved the verify gene's evolution. Two caveats attach. First, runs~7--10 gate the
loop on the true metric at \emph{held-out} scoring time as well as during training, so those
numbers --- including $+0.063$ and the detector-quality curve --- are \textbf{ceiling
results} conditional on having a perfect checker at inference, not deployable gains. Second,
at $n{=}3$ a sign test cannot fall below $p=0.25$ even when every seed agrees, so we now
report sign-test $p$ alongside every interval: $p=0.25$ for $+0.063$ and $p=0.50$ for the
actuator run.

\paragraph{The fix, and the reporting minimum we propose.}
The corrected harness (i) builds a \emph{component-specific} reflective dataset, so the verify
gene's records carry the draft it judged, the verdict it returned, and whether that verdict
was a false pass or a false alarm; (ii) passes an explicit component selector, defaulting to
rewriting all genes per proposal so a multi-gene arm is not starved; (iii) replaces arm~C
with a genome-blind constant draft scored by the real metric, which is flat without being
saturated, and re-measures arm~A instead of copying it; and (iv) records, for every arm, the
per-gene accepted-change count, the train aggregate, and the final gene text. Item~(iv) is
the part we would urge on others: \textbf{a componentized-optimization result should report,
per component, whether the optimizer's output actually differs from the seed.} Had we printed
that one number, this paper would have been a different paper eleven runs earlier. Because
the corrected wiring changes every reflection prompt, it changes every cassette key, so the
released cassettes cannot replay under it; they are retained and labelled superseded, and the
re-run under the corrected harness is future work.

\section{Threats to Validity}
\label{sec:threats}

We state these plainly. \textbf{Power, and what a $[0,0]$ interval actually means:}
runs 4--10 use $n{=}3$ seeds (runs 2--3 use 5). We previously argued that the intervals were
``unusually informative despite small $n$'' because the paired per-seed $\dd(D{-}B)$ has
near-zero variance, B and D often optimizing to byte-identical genomes. That argument was
circular and is withdrawn: byte-identical genomes mean the two arms issued the same calls and
the cassette served both the same completions, so the zero is an \emph{identity} rather than
a measurement, and an interval of $[0,0]$ is the sharpest available evidence that the
treatment was not applied (\S\ref{sec:defect}). We now report, per run, how many seeds were
such ties, alongside the interval and an exact sign test. Of the four fair null runs,
$3/5$ seeds are ties on run~3, $3/3$ on run~4, $2/2$ on the code replication and $2/3$ on the
model-swap replication. At $n{=}3$ the sign test cannot fall below $p=0.25$ even with
unanimous seeds, so the one positive effect, $\dd(\text{hi}-\text{lo})=+0.063$ with interval
$[+0.025,+0.090]$, rests on the magnitudes rather than on replication across seeds. There is
also no within-condition sampling variance anywhere in the study: the worker runs at the
provider's default temperature with no seed, one sample per task, and the cassette memoizes
by prompt, so every interval here describes task-draw and search-path variation only, never
measurement noise. \textbf{Single worker/reflection model:} results may be
model-dependent; \citet{huang2023cannot} find self-correction effects vary by
model. \textbf{Cross-lineage, not cross-provider:} run 6's critic
(\texttt{gpt-5.6-luna}) shares a vendor and training lineage with the worker
(\texttt{gpt-4.1-mini}), so it breaks self-\emph{grading} but may retain
self-preference bias; a true cross-provider critic (Gemini, Claude) is future work
(D3). \textbf{One task family:} verifiable instruction-following may be a
domain where the critique has little leverage; GEPA's own plateaus on
GPQA/MATH-500 suggest domain dependence. \textbf{Budget:} $\mathrm{mmc}\in[20,64]$
is small; a larger budget could in principle reveal a crossover. \textbf{Held-out
generalization:} the generator's parameter pools are small, so train/held-out
tasks can share parameter combinations; this affects all arms equally and does
not bias \dd, but softens the generalization claim. Section~\ref{sec:directions}
turns each threat into a concrete follow-up.

\paragraph{Which threats the released artifact can close.} The threats above are
not equally tractable, and the difference is not scientific interest but what each one
actually requires: more spend against the same harness, a credential we do not hold, or new
harness code. We separate them explicitly (Table~\ref{tab:closability}) so a reader can tell
which open questions are a budget line and which are engineering. The audit in
Section~\ref{sec:defect} reorders this list: the harness re-run now dominates everything
else, because until the loop-level treatment is actually applied the remaining threats are
threats to a measurement that did not happen.

We report one replication completed before the audit. A worker swap to
\texttt{gpt-4.1-nano} on the decisive multi-hard regime ($3$ seeds, $16$/$16$,
$\mathrm{mmc}=64$) gives $\dd(D-B) = -0.014$ (per-seed $+0.000, -0.041, +0.000$; corrected
interval $[-0.041,+0.000]$; two of three seeds are ties), with arm~B at
$0.70 \pm 0.03$ and arm~D at $0.69 \pm 0.03$ from a seed of $0.67$. It shows the measured
pattern is not an artifact of one worker model --- but it inherits the same harness, so like
the runs it replicates it is not a test of loop-level optimization. An $\mathrm{mmc}{=}160$
budget replication was designed, priced, started and then stopped mid-run once the audit
landed, having spent $\approx\!\$0.33$ of a projected $\$0.90$: continuing to buy data under
the superseded wiring was no longer worthwhile. Its partial cassette is not released, since a
recording with no corresponding result would be a misleading artifact.

\begin{table}[h]
\centering
\small
\begin{tabular}{llp{0.40\textwidth}}
\toprule
Threat & Closable as released? & What it requires \\
\midrule
\textbf{Treatment not applied} (\S\ref{sec:defect}) & \textbf{No --- re-run required} &
  The corrected harness exists; the runs do not. Blocks every row below \\
Single worker/reflection model & Partly, by spend & Worker swap done pre-audit
  ($\dd=-0.014$); needs redoing post-fix \\
Budget, $\mathrm{mmc}\in[20,64]$ & Yes, by spend & An $\mathrm{mmc}{=}160$ replication under
  the corrected harness \\
Power, $n{=}3$; no sampling variance & No --- design change & Repeated samples per task and a
  fixed temperature, neither of which the study has (D7) \\
Cross-lineage, not cross-provider & No --- credential & A Gemini or Anthropic key; pricing
  already in place (D3) \\
One task family (math) & No --- new harness & Parametric math generator plus numeric checker;
  instance-level visible/hidden split \\
Held-out generalization & No --- generator work & Enlarged parameter pools, then re-run and
  report the gap (D8) \\
\bottomrule
\end{tabular}
\caption{Threat closability against the released artifact and an OpenAI-only key. ``By
spend'' means the harness answers it with no new code. The first row was added after the
audit and supersedes the others in priority: it is not a threat to the interpretation of a
result but to the existence of one.}
\label{tab:closability}
\end{table}

\section{When Would Loop-Level Help? Directions}
\label{sec:directions}

The null is not ``componentized optimization is useless''; it is ``the critique
component of a self-refine loop is not a productive optimization target when it
is an unreliable, self-referential detector on shallow-verification tasks.'' The
following directions test the boundary; each changes one variable.

\emph{Done in this paper:} D2 (critique-as-detector, \S\ref{sec:analysis}), a
cross-lineage critic (run 6), the oracle ceiling (run 7), a fix-budget sweep
(run 8), and the feedback-actionability ablation (D2$'$, runs 9--10). What runs
7--10 changed: the priority is no longer ``get a better detector'' --- a perfect
one does not suffice --- but ``fix the fixer,'' and we now know the fixer is
\emph{partly} feedback-limited (actionability buys $\sim\!+0.06$) and \emph{partly}
generator-limited (the residual). The agenda, re-ordered accordingly:

\begin{description}
  \item[D1$'$ (\emph{fix}-side interventions).] The binding constraint is repair,
    not detection. Run 8 did the refine-budget half ($R{:}1\!\to\!3$, $+0.065$).
    \emph{Remaining:} a stronger worker on refinement only, and constraint-specific
    repair \emph{tools} (e.g.\ a deterministic word-count trimmer) that let the model
    \emph{act} on feedback --- the natural next step given the generator-limited
    residual runs 9--10 expose.
  \item[D2$'$ (feedback actionability) --- \textbf{done, runs 9--10}.] Holding
    detection at the perfect oracle, concrete \emph{measured} edits beat a terse
    verdict by $+0.063$ (seed $0.810\!\to\!0.873$), while naming the failed
    constraint alone adds nothing (lo $\approx$ mid). So the fix step is partly
    feedback-limited, the actionable content is the \emph{measurement} (not the
    diagnosis), and the effect is feedback shape, not GEPA ($\dd(D-B)=0.00$). Open
    follow-up: does an LLM \emph{rewriting} the measured edits into fluent prose add
    anything over the raw measurement?
  \item[D3 (true cross-family critic).] Repeat run 6 with a cross-\emph{provider}
    critic (\texttt{gemini-2.5-flash-lite} or \texttt{claude-haiku-4-5}; pre-priced,
    needs a key) at $\ge 8$ seeds, to test whether the $+0.01$ nudge is real.
  \item[D4 (budget \& worker sweeps).] $\mathrm{mmc}\in\{20,40,80,160\}$ and worker
    $\in$\{nano, mini, a 5-series\} to rule out under-budgeting and map
    model-dependence.
  \item[D5 (cross-domain) --- \textbf{done} (\S\ref{sec:bcb}, App.~\ref{app:mbpp-failed}).]
    We replicated the 4-arm design in a domain where the verify step catches real
    \emph{semantic} errors and runs as an executable tool, the regime
    \citet{shinn2023reflexion} find self-correction works. A first attempt on MBPP+
    \citep{liu2023evalplus} saturated its training split (Appendix~\ref{app:mbpp-failed});
    moving to frontier-hard BigCodeBench-Hard \citep{zhuo2024bigcodebench} restored headroom
    (seed $0.40$--$0.50$), and there GEPA accepted training-improving candidates for
    \emph{both} arms on one seed yet still delivered $\dd(D-B)=0.00$ on held-out
    (\S\ref{sec:bcb}). The instruction-following null thus recurs in a semantically-verified,
    executable-feedback domain: reliable, actionable verification is \emph{not} by itself
    sufficient to make the loop-level lever pay off. Aiming GEPA directly at the fixer's
    feedback (the \emph{actuator} experiment, \S\ref{sec:bcb}) did not rescue it either: a
    powered $3$-seed replication ($20$/$20$, $\mathrm{mmc}=50$) gave $\dd(D-B)=-0.10$
    (corrected CI $[-0.150,+0.000]$, which includes zero; sign test $p=0.50$), and the
    over-fitting explanation we attached to it does not survive the harness audit
    (\S\ref{sec:defect}). It therefore does \emph{not} close the budget/power threat. Math (a
    numeric checker) remains a natural next replication, and a re-run under the corrected
    harness now precedes all of these.
  \item[D7 (statistics) --- \emph{reinstated}.] The near-zero paired variance we once
    read as precision was mostly ties, i.e.\ the treatment not being applied
    ($\S\ref{sec:defect}$), and no interval in the study contains within-condition sampling
    variance at all. Repeated samples per task and more seeds are now load-bearing rather than
    optional. Retained from the original note: exact-resample intervals are reported, and
    them. A larger-$n$ pre-registered replication remains nice-to-have, not load-bearing.
  \item[D1$'$-tools (deterministic actuator).] Table~\ref{tab:residual} localizes the
    residual to \texttt{exact\_word\_count}; a deterministic count-fixer would likely
    recover most of it. Lower priority now that the decomposition has made the point
    (a counter fixes counting); its value is confirming the other constraints do not
    break when the count is enforced.
  \item[D8 (generalization).] Enlarge the generator's parameter space so
    train/held-out cannot near-duplicate, and report the train--held-out gap.
\end{description}

\section{Conclusion}
\label{sec:conclusion}

Three questions are easily conflated; our answers differ. \textbf{(i) Does the
self-refine loop help?} Only through \emph{external} signal --- a measurement tool
(seed $0.765\!\to\!0.837$) or more refine attempts ($0.748\!\to\!0.813$) --- never
through the model grading its own draft (the intrinsic critique is a near-random
detector, F$_1\!=\!0.11$). \textbf{(ii) Does GEPA lift this task?} No: its optimized
genomes gave no held-out gain and sometimes \emph{degraded} both arms (train-split
overfitting), at the budgets tested --- not a verdict on GEPA in general, which
reports gains elsewhere. \textbf{(iii) Does loop-level optimization beat prompt-only
(the thesis)?} \emph{We cannot say, and that is the paper's main finding.} Our harness did not
apply the loop-level treatment: the critique gene is byte-identical to the seed in $15/15$
winning genomes on the flagship run, because the reflector was shown the generator's outputs
rather than the critique's, and because a one-component-per-proposal selector starved the
two-gene arm on the gene that moves the metric ($\S\ref{sec:defect}$). The measured
$\dd(D{-}B)\approx 0$ is therefore consistent with the hypothesis being false \emph{and} with
its never having been tested; the $[0,0]$ intervals we once cited as precision are ties, which
is evidence for the latter.

Across ten pre-registered instruction-following experiments, a cross-domain replication on
frontier-hard code (\S\ref{sec:bcb}), an actuator variant, and a worker-model swap, we
measured $\dd(D{-}B)\approx 0$ every time --- and the audit in \S\ref{sec:defect} shows why
that consistency is not the evidence it appears to be: on the runs we inspected the critique
gene never changed, so the arms differed mainly in how much search each spent on the gene they
shared. A confident zero, replicated eleven times, from an experiment that did not run.

What we can still say divides cleanly. \textbf{About the loop:} the intrinsic critique is a
near-random failure detector ($\mathrm{F}_1 = 0.11$), which quantifies the practitioner
intuition that a model grading its own homework is too lenient
\citep{osmani2026loop,kamoi2024whencan}; \emph{given} a perfect checker at inference, a
perfect detector still does not maximize held-out score, so detection is not the binding
constraint there either; the repair signal's \emph{shape} is worth $+0.063$ (text) and
$+0.20$ (code) while naming the failed constraint adds nothing; and the residual concentrates
in one mechanical skill, exact counting, that no prompt rewrite fixed. \textbf{About
componentized optimization:} nothing yet --- and that is the point. The apparatus we built to
answer the question quietly answered a different one.

\textbf{About building such apparatus}, which is the contribution we did not plan. A
hill-climbing loop improves a harness by attributing trace evidence to components
\citep{langchain2026loopcraft}; if that attribution is wrong, the optimizer rewrites the wrong
thing and every downstream number stays plausible. Two defaults were enough to cause it here:
a reflective dataset shared across components, and a selector that rewrites one component per
proposal. Neither is exotic; one is a library default. The cheap defence is to report, per
component, whether the optimizer's output differs from the seed at all --- one integer per
gene, which we now emit and which would have caught this eleven runs earlier. We release the
corrected harness, every cassette and pre-registration including the superseded ones, and the
audit that produced this section, on the view that a study which invalidates its own headline
is more useful published than withdrawn.

\appendix

\section{A Failed First Attempt: MBPP+ Saturation, and What It Taught Us}
\label{app:mbpp-failed}

Our first attempt at the cross-domain code study (D5) used \textbf{MBPP+}
\citep{liu2023evalplus}: the base tests as the in-loop (visible) verifier and the
edge-case ``plus'' tests as the hidden scorer, with the same four arms and the same
\texttt{gpt-4.1-mini} worker as the main study. We report it in full because its
\emph{failure} is what motivated the design of the main cross-domain experiment
(\S\ref{sec:bcb}).

\paragraph{What happened.} The run completed and reproduces bit-for-bit from its
cassette, but it is \emph{uninformative} for the quantity of interest. Every arm
returned exactly the seed genome's held-out score on both seeds --- $0.80$ (seed 0)
and $0.90$ (seed 1) --- so $\dd(D{-}B)=0$ identically. This zero is not a null result:
GEPA's own log shows the cause --- \texttt{``All subsample scores perfect for parent
0. Skipping''} on every iteration. The seed loop, aided by executable base-test
feedback, already solved all six training tasks at a perfect score, so the optimizer
had \emph{no gradient} and proposed no candidates. The ``optimized'' genome for each
arm was therefore the seed genome, and $\dd(D{-}B)=0$ is a degenerate artifact of
optimization never running --- the same saturation failure mode as our earliest
IFBench pilots (runs 1--2).

\paragraph{Why MBPP+ saturates.} MBPP+ is entry-level Python, and a capable worker
with real unit-test feedback in the loop solves the easy training tasks outright.
Held-out headroom looked adequate ($0.80$--$0.90$), but the \emph{training} split had
none, which is what gates GEPA. The lesson: a cross-domain test of loop-level
optimization needs headroom on the \emph{training} tasks, not merely the held-out set.

\paragraph{Why we did not simply weaken the worker.} The obvious fix --- a smaller
worker (e.g.\ \texttt{gpt-4.1-nano}) --- manufactures headroom but destroys external
validity. A robust literature documents that both prompt-optimization and
self-refinement gains scale \emph{inversely} with model strength:
\citet{zhou2025bbopromptscale} measure black-box prompt-optimization lift falling from
$12\%$ (7B) to $5.9\%$ (72B) to $1.1\%$ (671B) on the same task, concluding that gains
measured on a small model do not transfer to a large one; and
\citet{zhang2024smallverifiers} find small models need a \emph{strong} external verifier
for self-correction to help at all. A positive $\dd(D{-}B)$ obtained by crippling the
worker would thus be an upper bound that attenuates toward zero at frontier scale ---
the opposite of a generalizing result. Conversely, the saturation we observed
generalizes \emph{upward}: a stronger model saturates MBPP+ even more, so ``loop-level
optimization has no headroom on easy code for a capable model'' only becomes more true
for 2026-era frontier models.

\paragraph{How this reshaped the study.} The fix is to move headroom to the
\emph{tasks}, not away from the model: use a benchmark that is hard even for frontier
models, so an affordable worker sits far from the ceiling in the same
capability-relative regime a frontier model occupies on frontier-hard work. We adopted
\textbf{BigCodeBench-Hard} \citep{zhuo2024bigcodebench}, a subset calibrated to stay hard
even for strong models. Its tasks use scientific libraries and mock real I/O, so scoring moved
into a pinned Docker sandbox (\texttt{-{}-network none}, memory/PID caps, single-thread
BLAS) --- a safety and reproducibility upgrade over the subprocess sandbox, not a
compromise of it. The anti-oracle split is reconstructed by partitioning each task's
\texttt{unittest} methods into a visible (in-loop) and a hidden (scored) half. That
study is reported in \S\ref{sec:bcb}; this appendix records the attempt that led to it,
and its cassette is released alongside the others.

\bibliographystyle{plainnat}
\bibliography{references}

\end{document}
